% docs/manual/latex/user/08_running_and_reproducibility.tex - chapter of the LaTeX edition of the book "user_manual" (SIMC++ documentation).
%   Converted once from docs/manual/markdown/user/08_running_and_reproducibility.md by tools/md2latex.py; from then on THIS FILE is the source of the
%   LaTeX edition and may be edited by hand. Re-running tools/md2latex.py --write-sources for this
%   chapter OVERWRITES the hand edits (see docs/README.md).
%   Compile: cmake --build <build> --target pdf, or tools/build_docs_pdf.py (figures, xelatex passes,
%   makeindex with docs/common/latex/simcpp.ist); by hand, see docs/README.md. XeLaTeX, class
%   docs/common/latex/simcpp.cls (packages: fontspec, tikz, xcolor, hyperref, listings, longtable, booktabs,
%   titlesec, titletoc, fancyhdr, makeidx, preview, ...); fonts TeX Gyre Pagella, TeX Gyre Heros, DejaVu Sans
%   Mono. The figures are built first from docs/{manual,translation}/figures/*.tex (preview.sty) into PDFs.
\chapter{\texorpdfstring{Running modes and reproducibility}{Running modes and reproducibility}}\label{user_manual-08u5frunningu5fandu5freproducibility}
\section{\texorpdfstring{Seeds}{Seeds}}\label{user_manual-08u5frunningu5fandu5freproducibility:81-seeds}
The seed comes from, in this order: the pipeline \code{"rng":\ \{"seed":\ N\}}, then \code{-{}-{}seed\ N}\simcidx{Command-line options!seed@\code{-{}-{}seed}}, then the input's \hyperref[param:randomu5fseed]{\code{random\_\allowbreak{}seed}}\simcidx{Parameters!random_seed@\code{random\_seed}}.
\code{-{}-{}seed}\simcidx{Command-line options!seed@\code{-{}-{}seed}} overrides \hyperref[param:randomu5fseed]{\code{random\_\allowbreak{}seed}}\simcidx{Parameters!random_seed@\code{random\_seed}}, and the \code{.\allowbreak{}hist} prints the overridden value. With the default generator:

\begin{itemize}
\item \code{random\_\allowbreak{}seed\ >\ 0}: \code{std::\allowbreak{}mt19937\_\allowbreak{}64} seeded with it;
\item \code{random\_\allowbreak{}seed\ =\ 0} or absent (and other values ≤ 0 except −1): the fixed seed 5489, so the run is still reproducible;
\item \code{random\_\allowbreak{}seed\ =\ −1}: the current time. The log line \code{\ Use\ random\ seed\ =\ .\allowbreak{}.\allowbreak{}.} shows the value.
\end{itemize}
A given input, seed, mode and number of threads always produce the same files.

\section{\texorpdfstring{Random-number generators: \code{mt64} and \code{ranlux}}{Random-number generators: mt64 and ranlux}}\label{user_manual-08u5frunningu5fandu5freproducibility:82-random-number-generators-mt64-and-ranlux}
\begin{simctable}{footnotesize}{3}{>{\raggedright\arraybackslash}p{\dimexpr 0.169\linewidth-2\tabcolsep\relax}>{\raggedright\arraybackslash}p{\dimexpr 0.282\linewidth-2\tabcolsep\relax}>{\raggedright\arraybackslash}p{\dimexpr 0.488\linewidth-2\tabcolsep\relax}}
%
\toprule\rowcolor{simchead}\simcth{generator} & \simcth{pipeline} & \simcth{properties} \\\midrule\endfirsthead
\multicolumn{3}{@{}l}{\simccontinued}\\\toprule\rowcolor{simchead}\simcth{generator} & \simcth{pipeline} & \simcth{properties} \\\midrule\endhead
\midrule\multicolumn{3}{r@{}}{\simccontinues}\\\endfoot
\bottomrule\endlastfoot
\code{mt64} (default) & \code{"rng":\ "mt64"} or \code{\{"type":\ "mt64",\ "seed":\ N\}} & std::\allowbreak{}mt19937\_\allowbreak{}64; fast; allows \code{-{}-{}threads}\simcidx{Command-line options!threads@\code{-{}-{}threads}} \\
\code{ranlux} & \code{"rng":\ \{"type":\ "ranlux"\}}, optionally \code{"seed":\ N} or \code{"state":\ [25\ integers]} & RANLUX at luxury level 3, the generator of the legacy SIMC, seeded from \hyperref[param:randomu5fseed]{\code{random\_\allowbreak{}seed}}\simcidx{Parameters!random_seed@\code{random\_seed}} the same way. About 25 \% slower. \code{"state"} takes the 25 integers of a saved generator state (a legacy \code{*\_\allowbreak{}start\_\allowbreak{}random\_\allowbreak{}state.\allowbreak{}dat} file) \\
\end{simctable}
\section{\texorpdfstring{Default and legacy-compatible mode}{Default and legacy-compatible mode}}\label{user_manual-08u5frunningu5fandu5freproducibility:83-default-and-legacy-compatible-mode}
SIMC++ is a translation of the legacy Fortran SIMC. The legacy program has a number of known bugs (for example, the cut
that decides whether an event counts in the \code{.\allowbreak{}hist} yield uses the electron delta limit for the hadron). \textbf{The default
mode corrects them. The legacy-compatible mode, \code{-{}-{}fortran-{}compat}\simcidx{Command-line options!fortran-compat@\code{-{}-{}fortran-{}compat}}, reproduces the legacy behaviour exactly}, bugs
included:

\begin{simctable}{small}{3}{>{\raggedright\arraybackslash}p{\dimexpr 0.387\linewidth-2\tabcolsep\relax}>{\raggedright\arraybackslash}p{\dimexpr 0.115\linewidth-2\tabcolsep\relax}>{\raggedright\arraybackslash}p{\dimexpr 0.430\linewidth-2\tabcolsep\relax}}
%
\toprule\rowcolor{simchead}\simcth{} & \simcth{default} & \simcth{\code{-{}-{}fortran-{}compat}} \\\midrule\endfirsthead
\multicolumn{3}{@{}l}{\simccontinued}\\\toprule\rowcolor{simchead}\simcth{} & \simcth{default} & \simcth{\code{-{}-{}fortran-{}compat}} \\\midrule\endhead
\midrule\multicolumn{3}{r@{}}{\simccontinues}\\\endfoot
\bottomrule\endlastfoot
known bugs of the legacy SIMC & corrected & kept \\
with \code{ranlux}: identical output files to a legacy SIMC run of the same input & no & yes \\
\code{-{}-{}threads\ N\ >\ 1}\simcidx{Command-line options!threads@\code{-{}-{}threads}} & yes & refused \\
use it for & physics & reproducing or checking an existing legacy SIMC result \\
\end{simctable}
In a JSON input: \code{"run":\ \{"fortran\_\allowbreak{}compat":\ true\}}. The log says which mode is active
(\code{\ SIMC++:\ fortran\_\allowbreak{}compat\ =\ true\ (\allowbreak{}exact\ Fortran\ behaviour,\ known\ Fortran\ bugs\ kept)}).

\begin{simclegacy}
what each correction changes in the results, and which run types are affected, is listed in
\code{docs/\allowbreak{}translation/\allowbreak{}markdown/\allowbreak{}user\_\allowbreak{}differences.\allowbreak{}md}; how to compare SIMC++ with legacy SIMC runs statistically is in
\code{docs/\allowbreak{}translation/\allowbreak{}markdown/\allowbreak{}user\_\allowbreak{}comparing\_\allowbreak{}with\_\allowbreak{}fortran.\allowbreak{}md}.
\end{simclegacy}
\begin{figure}[htbp]
\simcfigure{compat_switch}
\caption[The two behaviour modes]{The two behaviour modes: what each Category of finding does by default and with -{}-fortran-compat,\allowbreak{} and how the switch reaches the code}\label{fig:compatu5fswitch}
\end{figure}
\section{\texorpdfstring{Reproducing legacy SIMC results exactly}{Reproducing legacy SIMC results exactly}}\label{user_manual-08u5frunningu5fandu5freproducibility:84-reproducing-legacy-simc-results-exactly}
Run the same \code{.\allowbreak{}inp} with \code{ranlux}, \code{-{}-{}fortran-{}compat}\simcidx{Command-line options!fortran-compat@\code{-{}-{}fortran-{}compat}} and one thread. The \code{.\allowbreak{}hist} (apart from its two time lines), the
\code{.\allowbreak{}geni}, the \code{.\allowbreak{}gen} and every entry of the tree are then identical to those of the legacy program. Example with a stored
legacy run (\code{\$SIMC\_\allowbreak{}REFERENCE/\allowbreak{}runs/\allowbreak{}.\allowbreak{}.\allowbreak{}.} are the project's reference runs; 1000 events, with an \hyperref[param:extrau5fdbaseu5ffile]{\code{extra\_\allowbreak{}dbase\_\allowbreak{}file}}\simcidx{Parameters!extra_dbase_file@\code{extra\_dbase\_file}}):

\begin{simcshell}[basicstyle=\ttfamily\tiny]
echo '{"pipeline": {"rng": {"type": "ranlux"}}}' > ranlux.json
R=$SIMC_REFERENCE/runs/eep_h_test_pctnames            # the legacy run: .inp, extra file, .hist, .geni, .gen, .root
simc++ $R/eep_h_test_pctnames.inp -o exact --fortran-compat --pipeline ranlux.json > exact.log
diff <(sed 2,3d exact/eep_h_test_pctnames.hist) <(sed 2,3d $R/eep_h_test_pctnames.hist) && echo ".hist: identical apart from the time lines"
cmp exact/eep_h_test_pctnames.geni $R/eep_h_test_pctnames.geni && cmp exact/eep_h_test_pctnames.gen $R/eep_h_test_pctnames.gen && echo ".geni, .gen: identical"
simc_root_identity exact/eep_h_test_pctnames.root $R/eep_h_test_pctnames.root | tail -1
\end{simcshell}
Output:

\begin{simcplain}[basicstyle=\ttfamily\tiny]
.hist: identical apart from the time lines
.geni, .gen: identical
exact/eep_h_test_pctnames.root vs .../reference/runs/eep_h_test_pctnames/eep_h_test_pctnames.root: IDENTICAL
\end{simcplain}
Notes:

\begin{itemize}
\item Use the same input, including its \hyperref[param:randomu5fseed]{\code{random\_\allowbreak{}seed}}\simcidx{Parameters!random_seed@\code{random\_seed}} (none = the RANLUX default seed). Bit-for-bit identity also needs the
same C math library (glibc \code{libm}) as the machine that made the legacy run.
\item Without \code{-{}-{}fortran-{}compat}\simcidx{Command-line options!fortran-compat@\code{-{}-{}fortran-{}compat}}, the \code{.\allowbreak{}hist} differs in the corrected lines (here two lines where the legacy program writes
integers as \code{T}/\code{F}).
\item This holds for every example input distributed with SIMC that the legacy program can run (CTest label
\code{regression\_\allowbreak{}shipped}).
\end{itemize}
\section{\texorpdfstring{Threads: \code{-{}-{}threads\ N}}{Threads: -{}-threads N}}\label{user_manual-08u5frunningu5fandu5freproducibility:85-threads---threads-n}
\code{-{}-{}threads\ N}\simcidx{Command-line options!threads@\code{-{}-{}threads}} (or \code{"pipeline":\ \{"threads":\ N\}}; the option wins) runs the event loop on N threads
(\code{docs/\allowbreak{}manual/\allowbreak{}markdown/\allowbreak{}reference/\allowbreak{}PERFORMANCE.\allowbreak{}md} §3):

\begin{itemize}
\item each thread has its own copy of every stage and its own mt64 random stream, derived from (seed, thread number);
\item the events are split: thread k makes ngen/N events (+1 for k < ngen mod N), so the total is exactly ngen (successes
for ngen > 0, tries for ngen < 0);
\item the outputs are merged in a fixed thread order: the tree entries thread after thread, the histograms and counters
summed, the normalisation computed from the merged totals;
\item \textbf{a given seed and N always give the same files}, but they are statistically equivalent to, not equal to, the
one-thread run, or a run with another N. \code{-{}-{}threads\ 1}\simcidx{Command-line options!threads@\code{-{}-{}threads}} is exactly the sequential run;
\item the summary JSON records \code{"run":\ \{"threads":\ N\}};
\item the stage warnings of the threads may be interleaved in the log (\code{-{}-{}quiet}\simcidx{Command-line options!quiet@\code{-{}-{}quiet}} silences them).
\end{itemize}
\begin{simcshell}[basicstyle=\ttfamily\tiny]
$ simc++ run.json -o out_t4 --ngen 5000 --threads 4 > out_t4.log; grep -E "Ntried|Ncontribute" out_t4/run.hist
            Ntried         =     102516
            Ncontribute    =       5000
$ simc++ run.json -o out_t4b --ngen 5000 --threads 4 > /dev/null && simc_root_identity out_t4/run.root out_t4b/run.root | tail -1
out_t4/run.root vs out_t4b/run.root: IDENTICAL
\end{simcshell}
\begin{simcwarning}
\code{-{}-{}threads\ N\ >\ 1}\simcidx{Command-line options!threads@\code{-{}-{}threads}} does not work with \code{-{}-{}fortran-{}compat}\simcidx{Command-line options!fortran-compat@\code{-{}-{}fortran-{}compat}} or RANLUX: the exact reproduction of a legacy
run needs one random stream consumed in the legacy order, and N threads use N streams. Both combinations are refused
before the run starts (exit status 2):
\end{simcwarning}
\begin{simcshell}[basicstyle=\ttfamily\tiny]
$ simc++ run.json -o out_compat --ngen 5000 --threads 4 --fortran-compat
simc++: --threads 4 (more than one event-loop thread) cannot be used with --fortran-compat: the exact Fortran reproduction needs the single random stream in the Fortran order (run with --threads 1)
\end{simcshell}
Statistical comparisons of 8-thread and one-thread runs agree (\code{docs/\allowbreak{}manual/\allowbreak{}markdown/\allowbreak{}reference/\allowbreak{}PERFORMANCE.\allowbreak{}md} §3).

\begin{figure}[htbp]
\simcfigure{parallel_mode}
\caption[The parallel event loop]{The parallel event loop: the master builds N workers with their own stages and random streams, the workers run independently,\allowbreak{} and the results are merged in a fixed order}\label{fig:parallelu5fmode}
\end{figure}
\section{\texorpdfstring{Other run options}{Other run options}}\label{user_manual-08u5frunningu5fandu5freproducibility:86-other-run-options}
\begin{itemize}
\item \code{-{}-{}ngen\ N}\simcidx{Command-line options!ngen@\code{-{}-{}ngen}} overrides \hyperref[param:ngen]{\code{ngen}}\simcidx{Parameters!ngen@\code{ngen}} (use a small N for a test; a negative N gives |N| tries, which always ends).
\item \code{-{}-{}quiet}\simcidx{Command-line options!quiet@\code{-{}-{}quiet}} silences the stage messages (the progress lines stay).
\item \code{-{}-{}print-{}pipeline}\simcidx{Command-line options!print-pipeline@\code{-{}-{}print-{}pipeline}} prints the plug-ins chosen and continues with the run.
\item \code{-{}-{}run-{}name}\simcidx{Command-line options!run-name@\code{-{}-{}run-{}name}}, \code{-{}o}: output naming (\hyperref[user_manual-04u5finputs:422-run]{4.2.2}).
\end{itemize}
